\chapter{2018 Nobel Prize in Economics}
\textbf{Laureates: }William Nordhaus (USA), Paul Romer (USA)

\section*{Reason for Award}
Integrating climate change and technological innovation into long-run macroeconomic analysis

\section{What They Did}
William Nordhaus developed integrated assessment models that combine economic growth with climate science, showing how climate change affects economic welfare and how economic policies affect climate outcomes. His DICE model quantified the costs of climate change and the benefits of carbon pricing, demonstrating that optimal policy requires putting a price on carbon emissions. Paul Romer developed endogenous growth theory, showing that technological progress is not exogenous but results from intentional investment in research and development. His models demonstrated that ideas are non-rivalrous goods that generate increasing returns, explaining why knowledge accumulation drives long-run growth. Together, they provided frameworks for understanding how innovation and climate policy affect economic development.

\section{Impact on Economic Thinking}
Nordhaus and Romer transformed macroeconomic analysis by incorporating climate change and innovation as central phenomena. Nordhaus showed that climate change is not an external issue but a fundamental economic problem requiring pricing mechanisms. His work influenced the Paris Agreement and carbon pricing debates worldwide. Romer's endogenous growth theory explained why some countries grow faster than others and how intellectual property rights affect innovation incentives. This influenced patent policy, R\&D subsidies, and education policy. Their combined work showed that long-run growth depends on both technological progress and environmental sustainability. The integration of climate and growth models created a new research program addressing the most pressing challenges of the twenty-first century.

\section{Modern Perspectives Today}
Nordhaus and Romer's frameworks are essential for addressing contemporary challenges. Climate economists use integrated assessment models to evaluate carbon pricing, adaptation strategies, and climate justice. The economics of climate change has become a major research area with policy implications worldwide. Romer's endogenous growth theory informs debates about innovation policy, intellectual property, and the knowledge economy. Research on green growth, sustainable development, and the circular economy builds on their integrated approach. The COVID-19 pandemic highlighted the importance of innovation for public health, extending Romer's insights about ideas as public goods. As economies face the dual challenges of climate change and technological transformation, Nordhaus and Romer's frameworks provide essential analytical tools.
